\subsection{Punctured codon cube at the standard-code readout layer}\label{sec:bioreality-h0-r-punctured-cube}
\origin{ai}

\paragraph{Finite object.}
This record studies a finite subset of the six-bit codon cube $Q_6$ at the level where codons are read only as standard-code labels.  The observed set $R$ is the thirteen-codon row
$$
\begin{aligned}
R=\{&\mathrm{AAA},\mathrm{AGA},\mathrm{AGG},\mathrm{AUA},\mathrm{CUA},\mathrm{CUC},\mathrm{CUG},\mathrm{CUU},\\
&\mathrm{UAA},\mathrm{UAG},\mathrm{UCA},\mathrm{UGA},\mathrm{UUA}\}.
\end{aligned}
$$
The cardinality check reports $|R|=13$, and the same thirteen names appear in the direct punctured-cube equality check.  The measured face data are
$$
\begin{aligned}
f(R)=(f_0,f_1,f_2)=(13,16,4),
\end{aligned}
$$
so the local tile has thirteen vertices, sixteen edges, and four square faces inside the codon-coordinate graph.  This is a finite row statement about a labeled coordinate pattern.  It is not a statement that the genetic code is generated by the geometry of $Q_6$, and it is not a statement about translation dynamics, folding, physical admissibility, biological function, mechanism, or universality.

\paragraph{Readout contact.}
The only biological contact used here is the curated standard-code table assigning codons to amino-acid or stop labels.  That contact is adequate for code-layer readback: it says which standard labels are attached to the thirteen codons above, and therefore permits finite checks about local label splits on $R$.  The contact does not observe ribosomal realization, molecular structure, thermodynamic admissibility, organismal function, or a general law over all biological coding systems.  Consequently every claim in this subsection is confined to the code-read layer.  Any stronger assertion would require an additional contact whose measurements live at the stronger layer.

\paragraph{Spectral row.}
The finite adjacency calculation supplies a numerical spectrum summary for the same coordinate object.  The reported values are $\lambda_M=0.6752479308830213$, $\lambda_R=0.4758352843197002$, and $\mu=3.051487585298128$.  The associated radial polynomial residual is $-6.224354365258478\cdot 10^{-12}$, which is also the reported cubic-polynomial residual.  These numbers are used only as finite coordinate diagnostics.  They support the claim that the punctured-cube row has a reproducible graph signature under the implemented calculation; they do not convert a coordinate signature into a biochemical mechanism.

\paragraph{Shape rarity checks.}
Against the recorded null model with seed $42$ and $N=2000$ samples, the exact shape of the observed object appeared once.  The exact-shape estimate is therefore $p=0.0005$.  In the same experiment, there were $25$ samples with Jaccard score at least $0.9$, giving $p=0.0125$ for the high-overlap event.  The run required $57377$ distinct-sample attempts and recorded $55377$ duplicate retries.  A second null model preserving amino-acid class information reports $N=2000$, zero exact shape matches, observed $M$-size $20$, and $7$ samples in the lower tail at or below the observed $M$-size.  Thus the class-preserving exact-shape estimate is $0.0$ in this finite sample, and the lower-tail estimate is $0.0035$.  These are empirical finite-row measurements under specified sampling procedures.  They should be read as evidence that the row is non-generic relative to those procedures, not as evidence for a universal biological necessity.

\paragraph{Predictive stress tests.}
The leave-one-code-out experiment evaluated $23$ tables and $63$ total targets.  Its baseline hit rate was $0.19821739130434785$.  The observed rate above baseline was $0.9841269841269841$, the rate inside $M-T$ was $0.9365079365079365$, and the rate inside $M-T$ plus the boundary was $0.9841269841269841$.  The leave-one-codon-out experiment reports recovery rate $1.0$, predictable count $8$, closure-coverage-high count $8$, boundary count $5$, and anchor count $0$.  These stress tests say that, within the finite procedure used here, omitted code or codon entries are usually recovered by the local row plus boundary information.  They do not say that translation itself has been reconstructed, nor that a protein sequence has acquired structure or function.

\paragraph{Family and module accounting.}
The family-resolved leave-one-code-out row separates five families.  The hit-rate-above-baseline values are $0.8333333333333334$ for alternative yeast, $1.0$ for bacterial or plastid, $1.0$ for candidate other, $1.0$ for mitochondrial, and $1.0$ for nuclear ciliate.  The module classification is complete on the thirteen-codon row: $R$-size is $13$, assigned count is $13$, and the lists of missing, extra, and unassigned codons are all empty.  The six recorded modules are
$$
\begin{aligned}
\mathrm{agr}&=\{\mathrm{AGA},\mathrm{AGG}\},\\
\mathrm{ar\_lys}&=\{\mathrm{AAA}\},\\
\mathrm{cun}&=\{\mathrm{CUA},\mathrm{CUC},\mathrm{CUG},\mathrm{CUU}\},\\
\mathrm{ile\_met}&=\{\mathrm{AUA}\},\\
\mathrm{ssn}&=\{\mathrm{UCA},\mathrm{UUA}\},\\
\mathrm{stop\_trp}&=\{\mathrm{UAA},\mathrm{UAG},\mathrm{UGA}\}.
\end{aligned}
$$
Every listed module has either predictable or boundary status in the recorded classification.  This closes the finite module bookkeeping for $R$: every codon in the row is assigned, and no codon outside the row is accidentally imported into the module list.

\paragraph{Interpretation.}
The result is a bounded code-read claim: a thirteen-codon subset of the six-bit codon cube agrees with a punctured local tile having face vector $(13,16,4)$, nontrivial spectral diagnostics, rare-shape behavior under two null procedures, and complete finite module accounting.  The standard-code table supplies the labels needed to read the row, while the coordinate construction supplies the tile geometry.  The two ingredients together support local codon-label analysis.  They do not support promotion to claims about ribosomal translation, folded protein structure, physical possibility, biochemical function, global mechanism, or universality.  Those stronger layers remain unclaimed unless separate measurements are attached to them.